\documentclass[11pt]{article}
\usepackage[margin=1in]{geometry}
\usepackage{amsmath,amssymb,amsthm,mathtools}
\usepackage[hidelinks]{hyperref}
\newtheorem{theorem}{Theorem}
\newtheorem{lemma}[theorem]{Lemma}
\newtheorem{proposition}[theorem]{Proposition}
\newtheorem{corollary}[theorem]{Corollary}
\theoremstyle{remark}\newtheorem{remark}[theorem]{Remark}
\newcommand{\an}{\mathrm{an}}
\newcommand{\Gm}{\mathbb G_m}
\newcommand{\Z}{\mathbb Z}
\newcommand{\R}{\mathbb R}
\newcommand{\M}{\mathcal M}
\newcommand{\OO}{\mathcal O}
\newcommand{\Pic}{\operatorname{Pic}}
\newcommand{\val}{\operatorname{val}}
\newcommand{\Etale}{\'etale}
\title{A strict proper analytic surface with constant meromorphic field:\\
an explicit torus and a nonalgebraization proof}
\author{Research working paper for Problem P16.7}
\date{8 October 2026; continuation 374; independent review pending}
\begin{document}
\maketitle
\begin{abstract}
Let $K=\mathbb C((s))$, $|s|=e^{-1/8}$, and $t=s^8$.
The quotient of $(\Gm^2_K)^{\an}$ by the periods $(t,2)$ and $(3,t)$
is a good strict smooth proper geometrically connected analytic surface.
We prove that its global meromorphic functions are constants. The argument
uses analytic triviality of line bundles on the covering torus, supplies
compatible trivializations on its entire exhaustion, and excludes every
nonzero indefinite theta multiplier by an explicit coefficient recurrence.
It excludes algebraization by a separated finite type algebraic space,
including a proper scheme, without first descending the analytic group law
or assuming that a hypothetical algebraization is projective.
This proves a geometric ingredient of the proposed example. It does not
construct the unrestricted intrinsic finite energy cone or solve the full
Monge--Ampere problem P16.7. The paper is an unaccepted working proof.
\end{abstract}

\section{The quotient and the precise claim}
Set $T=(\Gm^2_K)^{\an}$, with coordinates $z_1,z_2$, and put
\[
 \lambda_1=(t,2),\qquad\lambda_2=(3,t),\qquad
 \lambda(n)=(t^{n_1}3^{n_2},2^{n_1}t^{n_2})\quad(n\in\Z^2).
\]
Write $\Lambda=\lambda(\Z^2)$ and
$x=(-\log|z_1|,-\log|z_2|)$. Then $x(\lambda(n)z)=x(z)+n$.
In particular the lattice action is faithful and free. The calculation
uses actual equalities in characteristic zero, rather than just their
absolute values: $2^a=3^b$ for integers $a,b$ forces $a=b=0$.

\begin{theorem}\label{main}
The ordinary analytic quotient $X=T/\Lambda$ exists and is good, strict,
smooth, separated, proper, geometrically connected and pure of dimension two.
For every complete isometric extension $L/K$,
\[
 \Gamma(X_L,\M_{X_L})=L.
\]
Every analytic line bundle on $X_L$ with a nonzero global analytic section
is trivial, and that section has no zeros. Consequently $X_L$ has no
nonzero effective analytic Cartier divisor. There is no separated finite
type $L$-algebraic space $Y$ with $Y^{\an}\simeq X_L$; in particular there
is no proper scheme or proper algebraic space algebraization.
\end{theorem}

The analytic factoriality input used below is the following classical
fact, recorded in van der Put \cite[3.4.8]{vdP}: if $A$ is a regular
factorial affinoid algebra, then $A\langle u\rangle$ and
$A\langle u,u^{-1}\rangle$ are factorial. It is attributed there to Gruson
\cite{Gruson}. Finite affinoid coherent descent, regularity of smooth
polyannuli, Laurent expansions and multiplicativity of their Gauss norms
are also analytic foundations of the proof. These are cited mathematical
inputs; executing finite computations is not a proof of these inputs.
The complete Gruson unit freshly checked here is IV.5, Theorem 1,
for products of discs and unit circles. Its statement does not itself
give the general regular-affinoid version of van der Put 3.4.8 used
above, especially for a wide-polyannulus base. That additional
generality is an explicit source-scope verification gate for independent
acceptance. The argument below uses the printed classical assertion
as a cited foundation; it does not claim to have derived its full
generality from Gruson IV.5.
Sections \ref{units}--\ref{theta} give the argument from them, including
the exhaustion compatibility that finite polyannulus triviality alone
would not establish.

\begin{proof}[Construction of $X$ and its geometric properties]
If $B\subset\R^2$ is an open rectangle of side lengths less than one,
$U_B=x^{-1}(B)$ is ordinary open in $T$ and its distinct lattice translates
are disjoint. Give its image in the quotient its transported analytic
structure. On an overlap of two such images the lifts differ by the
unique locally constant lattice element. Its action is multiplication
by an invertible $K$-point, hence an analytic isomorphism. These maps
satisfy the cocycle relation and give an ordinary analytic atlas for $X$.

The action is properly discontinuous in a stronger compact sense:
for compact $C,D\subset T$, an element with $\lambda(n)C\cap D\ne\varnothing$
has $n\in x(D)-x(C)$, so there are only finitely many such $n$.
The graphs of these finitely many translations are closed. The action
map $\Lambda\times T\to T\times T$ is therefore proper, and the quotient
is Hausdorff. Its analytic diagonal is locally the finite disjoint union
of these closed graphs, which proves separatedness.

The strict compact polyannulus $x^{-1}([0,1]^2)$ maps onto $X$.
Every vector in $\R^2$ can be moved into $[0,1]^2$ by an integer vector;
the corresponding translation is an actual period. Thus $X$ is compact.
For goodness choose chart rectangles of side length $3/4$ with centers
in $(\tfrac18\Z)^2$. Given a point, take a center whose two coordinates
differ from its tropical coordinates by at most $1/16$. The closed
rectangle of side length $1/2$ with that center contains the point in
its ordinary interior and lies inside the chart rectangle. Its endpoints
are logarithms of elements of $|K^\times|$, so its inverse image is a
strict affinoid neighborhood. This proves goodness and strictness. Each
ordinary chart is an open subspace of the smooth boundaryless two-torus.
Hence $X$ is smooth and boundaryless, pure of dimension two; compactness,
separatedness and absence of boundary imply properness over $K$.

All charts and their overlap translations base change to any complete
$L/K$, and construct $X_L=T_L/\Lambda$ with the same properties. The
increasing polyannuli $|s|^m\le |z_i|\le |s|^{-m}$ covering $T_L$ have
integral affinoid algebras, so are connected, with nonempty nested
intersections. Therefore $T_L$ and its quotient are connected. Since
the local analytic rings are regular, irreducible components cannot
meet and are locally open. This also gives geometric integrality.
\end{proof}

\section{Units on a polyannulus and on the covering torus}\label{units}
The next estimates hold over an arbitrary complete nontrivially valued
field $L$, with $0<|\pi|<1$. In dimension $d$, let
\[
 P_m=\{ |\pi|^m\le |z_i|\le |\pi|^{-m}\ (1\le i\le d)\},
 \qquad A_m=\OO(P_m),\qquad a=-\log|\pi|>0.
\]
Every $f\in A_m$ has a unique Laurent expansion
$f=\sum_{\beta\in\Z^d}b_\beta z^\beta$ converging in the polyannulus norm.
For $y\in[-ma,ma]^d$, the multiplicative Gauss norm at the point with
$\log|z_i|=y_i$ gives
\[
 N_f(y)=\log\|f\|_y
       =\max_\beta(\log|b_\beta|+\beta\cdot y).
\]

\begin{lemma}\label{unitestimate}
If $f\in A_m^\times$, there are $b\in L^\times$, $\alpha\in\Z^d$ and
$g\in A_m$ such that $f=bz^\alpha(1+g)$, the coefficient of $z^0$
in $g$ is zero, and, for every integer $0\le j<m$,
\[
 \|g\|_{P_j}\le |\pi|^{m-j}.
\]
On $(\Gm_L^d)^{\an}$ every global analytic unit is exactly a scalar
times a Laurent monomial.
\end{lemma}
\begin{proof}
Both $N_f$ and $N_{f^{-1}}$ are convex, and multiplicativity gives
$N_f+N_{f^{-1}}=0$. Thus $N_f$ is affine on the box. On its interior
some term attaining the maximum has integral slope, so
$N_f(y)=\log|b_\alpha|+\alpha\cdot y$ with $b_\alpha\ne0$.
For every other nonzero coefficient, comparison at all vertices gives
\[
 |b_\beta/b_\alpha|\exp(ma\|\beta-\alpha\|_1)\le1.
\]
Take $b=b_\alpha$ and divide by $bz^\alpha$. On the smaller box this
bounds every nonconstant term by
$\exp(-(m-j)a\|\beta-\alpha\|_1)\le|\pi|^{m-j}$. The Gauss norm
formula gives the claimed estimate. The reasoning does not assert that
$g$ is small on the outer boundary $P_m$.

For a global unit apply the same convexity argument on all of $\R^d$.
The resulting affine function has one integral slope $\alpha$.
Any additional coefficient with exponent $\beta\ne\alpha$ would exceed
this affine function along a ray in direction $\beta-\alpha$.
Hence all such coefficients vanish and $f=bz^\alpha$.
\end{proof}

\section{Complete triviality of line bundles on the torus}
\begin{lemma}\label{finitepic}
Every $P_m$ is regular and factorial, and $\Pic(P_m)=0$.
\end{lemma}
\begin{proof}
We give the induction underlying \cite[3.4.9]{vdP}. The field is the
dimension zero case. Suppose that $B$ is the regular factorial algebra
of the first $d-1$ wide annuli. For a line bundle on their product with
$|\pi|^m\le|z|\le|\pi|^{-m}$, restrict to the boundary $|z|=|\pi|^m$.
Its algebra is $B\langle u,u^{-1}\rangle$ after scaling $z$, and is
regular factorial by the stated factoriality input. The boundary line
bundle is trivial. Use one trivialization to glue the given line bundle
with the trivial bundle on $|z|\le|\pi|^m$. The annulus and this disc
form a finite admissible affinoid cover of $|z|\le|\pi|^{-m}$.
Coherent descent produces an invertible sheaf on that larger disc.
Its algebra is $B\langle v\rangle$, also regular factorial, so the
extension is trivial. Its restriction to the annulus is trivial.
Thus the next wide polyannulus has trivial Picard group; its regularity
and the identification of divisor class group with Picard group make
its algebra factorial. This completes the induction.
\end{proof}

\begin{proposition}\label{globalpic}
For $T_L=(\Gm_L^d)^{\an}$ one has $\Pic(T_L)=0$.
\end{proposition}
\begin{proof}
Let $\mathscr L$ be an analytic line bundle. Choose frames $e_m$ on
$P_m$ for $m\ge1$, which is possible by Lemma \ref{finitepic}.
When a frame $e_m$ is already chosen, a tentative frame on $P_{m+1}$
restricts to $u_me_m$ for a unit $u_m\in A_m^\times$.
By Lemma \ref{unitestimate}, $u_m=b_mz^{\alpha_m}(1+g_m)$.
Rescale the tentative frame by the global monomial $b_m^{-1}z^{-\alpha_m}$.
The resulting $e_{m+1}$ satisfies
\[
 e_{m+1}|_{P_m}=(1+g_m)e_m,\qquad
 \|g_m\|_{P_j}\le |\pi|^{m-j}\quad(j<m).
\]
For fixed $j$, the product of the factors with $m>j$ converges in
$A_j$ to a unit: their differences from one tend to zero, and the
inverse products converge by the same bound. The first factor with
$m=j$ is a unit already, even though it need not be close to one.
Define
\[
 E_j=e_j\prod_{m=j}^{\infty}(1+g_m)|_{P_j}.
\]
The transition relation gives $E_{j+1}|_{P_j}=E_j$.
These frames glue on the admissible exhaustion to a global nowhere
vanishing frame of $\mathscr L$. This proves the global assertion,
rather than inferring it from independent finite-annulus triviality.
\end{proof}

\section{Periods force indefinite multipliers, and convergence excludes them}\label{theta}
Return to $d=2$ and the periods in Section 1, over any complete isometric
$L/K$. Pull an analytic line bundle on $X_L$ back to $T_L$ and choose
a global frame by Proposition \ref{globalpic}. A global section is
represented by a global analytic function $\theta$ with equations
\[
 \theta(\lambda z)=c_\lambda z^{h_\lambda}\theta(z),
 \qquad c_\lambda\in L^\times,\quad h_\lambda\in\Z^2.
\]
The multipliers are analytic units because they are transition
isomorphisms of the line bundle. By Lemma \ref{unitestimate} they have
exactly the displayed monomial form. Their cocycle law implies
\[
 h_{\lambda\lambda'}=h_\lambda+h_{\lambda'},\qquad
 \lambda^{h_{\lambda'}}=(\lambda')^{h_\lambda}.
\]
These assertions also hold for a common theta representation of a
meromorphic function, as shown in Section \ref{mero}.

\begin{lemma}\label{periodmatrix}
There is an integer $k$ such that
$h_{\lambda_1}=(0,k)$ and $h_{\lambda_2}=(k,0)$.
If $k\ne0$, no nonzero global analytic $\theta$ satisfies these equations.
If $k=0$, the vector space of such functions has dimension at most one
and every nonzero member is a Laurent monomial.
\end{lemma}
\begin{proof}
Write $h_{\lambda_1}=(A,B)$ and $h_{\lambda_2}=(C,D)$.
Symmetry gives $t^C2^D=3^At^B$.
The $s$-valuation gives $C=B$; the remaining constant equality
$2^D=3^A$ gives $A=D=0$. Put $k=B=C$.

Suppose $k\ne0$. Set $n=(1,1)$ if $k>0$, and $n=(1,-1)$ if $k<0$.
For $\lambda=\lambda(n)$ and $h=h_\lambda=(kn_2,kn_1)$,
\[
 h\ne0,\qquad Q=\val(\lambda^h)=2kn_1n_2=2|k|>0,
 \quad\val=-\log|\cdot|,\quad\val(t)=1.
\]
Take a nonzero coefficient $a_m$ in the everywhere convergent expansion
$\theta=\sum_{m\in\Z^2}a_mz^m$. The multiplier equation for this one
period gives, by uniqueness of Laurent coefficients,
\[
 a_{m+h}\lambda^{m+h}=c_\lambda a_m,
 \qquad
 a_{m+jh}=a_m c_\lambda^j
             \lambda^{-jm-hj(j+1)/2}\quad(j\ge0).
\]
Consequently
\[
 \val(a_{m+jh})=
 \val(a_m)+j\val(c_\lambda)-j\val(\lambda^m)
                 -\tfrac12j(j+1)Q\longrightarrow-\infty.
\]
But $\|m+jh\|\to\infty$, while convergence on the unit polyannulus
$|z_1|=|z_2|=1$ requires $a_{m+jh}\to0$. This contradiction excludes
every $k\ne0$, with the sign and convergence checked explicitly.

If $k=0$, a nonzero coefficient $a_m$ instead satisfies
$\lambda_i^m=c_{\lambda_i}$ for $i=1,2$. If also $a_{m'}\ne0$,
then $\lambda_i^{m-m'}=1$. Their valuations give $m_1=m'_1$
and $m_2=m'_2$. Hence there is at most one exponent, determined by
the two constants, and at most one independent function.
\end{proof}

It follows at once that a nonzero section of an analytic line bundle
on $X_L$ pulls back to a scalar monomial and has no zeros. It therefore
trivializes that line bundle downstairs. In particular a nonzero
effective Cartier divisor cannot exist.

\section{The full global meromorphic field}\label{mero}
Let $F$ be a nonzero global meromorphic function on $X_L$ and pull it
back to $T_L$. Smoothness makes the covering torus locally factorial.
Its pole divisor $P$ is therefore an effective analytic Cartier divisor.
By Proposition \ref{globalpic}, its ideal $\OO(-P)$ has a global frame;
its inclusion in $\OO$ is multiplication by a global analytic function
$\theta_0$ with divisor $P$. Then $\theta_1=F\theta_0$ is holomorphic:
there is no codimension one pole, and the normal affinoid rings equal
the intersection of their height one valuation rings. The pair
$\theta_0,\theta_1$ has no common local divisor.

Lattice invariance of $F$ implies
\[
 \theta_1(\lambda z)\theta_0(z)
       =\theta_1(z)\theta_0(\lambda z).
\]
At every prime divisor, local coprimeness of both pairs shows that
$\theta_0(\lambda z)/\theta_0(z)$ has valuation zero. Normality applied
to this quotient and its inverse makes it a global analytic unit.
Thus both $\theta_i$ have the same monomial multipliers. This is the
theta representation of \cite[4.3.1--4.3.4]{vdP}, with its divisor
argument and global triviality made explicit here. Lemma
\ref{periodmatrix} now forces both functions to be scalar multiples
of the same monomial. Hence $F\in L$, which proves
$\Gamma(X_L,\M)=L$. The zero function is already a constant.

\section{Excluding algebraic spaces, without projectivity}
Suppose that a separated finite type algebraic space $Y/L$ has
$Y^{\an}\simeq X_L$. The usual local faithfully flat analytification
and Jacobian criterion reflect reducedness, smoothness and dimension.
Connectedness and the regular local rings then make $Y$ integral of
dimension two. Alternatively these properties can be checked on an
\'etale affine atlas and at its closed points.

For such a Noetherian normal separated algebraic space, its generic
point lies in a dense open subscheme. Its rational function field
$L(Y)$ has transcendence degree two. The exact compatibility of rational
and meromorphic functions used here is Stacks tags
\href{https://stacks.math.columbia.edu/tag/0ENA}{0ENA} and
\href{https://stacks.math.columbia.edu/tag/0END}{0END}.
Normality gives only minimal associated points, finite type gives
finitely many generic points on each quasi-compact \Etale chart,
and the generic point is represented by its field inside the schematic
locus. These check all hypotheses of 0ENA; in particular no unrestricted
claim about arbitrary algebraic spaces or zero divisors is needed.

The following local argument checks the analytic passage. On an \Etale
affine atlas a rational
function is locally a quotient of regular functions by a nonzero
divisor. Analytification is flat and preserves these nonzero divisors.
The quotients agree on the analytified overlap and descend as
meromorphic functions. On a finite \Etale affinoid neighborhood this
is the faithfully flat equalizer for total fraction rings: for a
reduced affinoid $A$ and finite \Etale $A$-algebra $B$,
\[
 Q(B)=Q(A)\otimes_A B,\qquad
 Q(A)\longrightarrow Q(B)\rightrightarrows Q(B\otimes_A B).
\]
The equality follows by localizing at the nonzero divisors of $A$;
the remaining rings are products of finite separable fields. An
analytic \Etale chart is locally a finite \Etale chart followed by
an open immersion, so this check gives the required sheaf descent.
This is also the local total-fraction argument of Stacks tag
\href{https://stacks.math.columbia.edu/tag/0EN6}{0EN6}; the analytic
flatness and local finite \Etale description are standard analytic
foundations, not a consequence of an algebraic lemma alone.

The map $L(Y)\to\Gamma(Y^{\an},\M)$ is injective. To see this, restrict
to a nonempty affine open scheme $\operatorname{Spec}(A)\subset Y$.
A nonzero numerator $g\in A$ is nonzero at some closed point after
excluding its zero set and that of a denominator. The finite residue
field of this point carries the extended complete absolute value,
so it gives an analytic point where the quotient is nonzero.
The map fixes $L$. It would therefore embed a field of transcendence
degree two into $\Gamma(X_L,\M)=L$, a contradiction.
This completes the proof of Theorem \ref{main}.

\section{Source boundaries and the remaining original problem}
The complete printed units of van der Put used here are 3.1.1 and
3.1.3; 3.4.2--3.4.9; 4.3.1--4.3.4; and 4.4.1--4.4.3.
The whole statements and printed immediate proofs were read from the
preserved source and checked against PDF pages where OCR omits the
formulas. The published passage from finite annuli to the full torus
is replaced here by the convergent transition product in Proposition
\ref{globalpic}. The general theta dimension and polarization theorem
is not needed: our recurrence handles this lattice directly.
The source's 4.5 uses the expression ``algebraically independent'' for
a finite homogeneous family where its dimension argument needs linear
independence; that printed wording is not used as a theorem here.
Papikian \cite[4.2]{Papikian} confirms the multiplier setup, but its
Appell--Humbert and ampleness discussion defers essential proofs and
is not substituted for those proofs.

The applied current Stacks statements and complete immediate proofs
are 0ENA, 0END, 0EN6, 0BB1 and 086U. Noetherianity checks the finite
generic-point hypothesis 0BB1; separatedness ensures decency, and
086U supplies the dense schematic locus. Recursive proofs of the
analytic factoriality foundation, finite affinoid coherent descent,
faithfully flat analytification and local finite \Etale structure are
cited foundations. A separate primary-reading ledger records what
has actually been read; a retained hash is not a new physical reading.
The full source bodies and PDF pixels remain local research evidence.

This paper supplies a geometric nonalgebraization argument for the
specific quotient; it does not assert that every nonpolarizable
presentation is nonalgebraizable. It neither defines all intrinsic
singular psh weights nor proves full mass decreasing continuity,
global analytic obstacle orthogonality, finite energy duality for
arbitrary analytic measures, or unrestricted existence and uniqueness
for P16.7. Restricted bounded or separable results in continuation 373
are separate unaccepted results. The original P16.7 remains open and
the verified original-problem count is unchanged. Independent whole
proof and source reconstruction is required before accepting this
working paper; successful compilation would only check its typesetting.

\begin{thebibliography}{9}
\bibitem{vdP} M.~van der Put, \emph{Rigid analytic spaces}, Groupe
d'\'etude d'analyse ultram\'etrique 3 (1975--1976), no.~2,
expos\'e J7, J1--J20.
\url{https://www.numdam.org/item/GAU_1975-1976__3_2_A6_0/}.
\bibitem{Gruson} L.~Gruson, \emph{Fibr\'es vectoriels sur un polydisque
ultram\'etrique}, Ann. Sci. \'Ecole Norm. Sup. (4) 1 (1968), 45--89.
\url{https://www.numdam.org/item/ASENS_1968_4_1_1_45_0/}.
\bibitem{Papikian} M.~Papikian, \emph{Rigid analytic geometry and
abelian varieties}, Snowbird notes, 2004, Section 4.2.
\url{https://math.stanford.edu/~vakil/snowbird/mihranjun21.pdf}.
\bibitem{Stacks} The Stacks Project, tags 0ENA, 0END, 0EN6, 0BB1 and
086U, current primary pages read 8 October 2026.
\url{https://stacks.math.columbia.edu/}.
\end{thebibliography}
\end{document}
